\documentclass{beamer}

\usetheme{Berlin}
\usecolortheme{Orchid}
\useoutertheme{miniframes}
\setbeamertemplate{navigation symbols}{}

\usepackage[utf8]{inputenc}
\usepackage[T1]{fontenc}
\usepackage{amsmath}
\usepackage{amssymb}
\usepackage{booktabs}
\usepackage{graphicx}
\graphicspath{{../images/}}

\usepackage{xcolor}
\usepackage{listings}
\usepackage{hyperref}
\usepackage{tikz}
\usetikzlibrary{arrows.meta,positioning,shapes.geometric,calc}

\lstset{
  language=Java,
  basicstyle=\ttfamily\small,
  keywordstyle=\color{blue},
  commentstyle=\color{gray},
  stringstyle=\color{teal},
  showstringspaces=false,
  columns=fullflexible,
  breaklines=true
}

% Per-slide source line (references shown on the slide itself).
\newcommand{\slidesources}[1]{%
  \vfill
  {\tiny\color{gray}\raggedright\sloppy\parbox{\linewidth}{Sources: #1}\par}%
}

% Unit-wise source strings for this subject (used by slides/notes).
% Path is relative to the lecture `latex/` folder (the compilation CWD).
% OOPS with Java (CSEG1044) — unit-wise sources
%
% These are short, student-facing reference strings intended to be used with:
%   - \slidesources{...}  (in beamer slides)
%   - \notesources{...}   (in lecture notes)
%
% Style inspiration: other subjects in My-Subjects/ use semicolon-separated sources
% and include an "accessed" date for web documentation.

\newcommand{\OOPSUnitISources}{Schildt, \textit{Java: A Beginner's Guide} (10e), McGraw-Hill (Java fundamentals + OOP basics).; Oracle Java Tutorials: Getting Started + Language Basics + Classes and Objects (accessed \today).; Java SE API docs (e.g., \texttt{String}, \texttt{Scanner}) (accessed \today).}

\newcommand{\OOPSUnitIISources}{Schildt, \textit{Java: The Complete Reference} (12e), McGraw-Hill (classes, inheritance, packages, interfaces).; Bloch, \textit{Effective Java} (3e), Addison-Wesley, 2018 (design choices; interfaces vs abstract classes).; Oracle Java Tutorials: Classes and Objects + Interfaces + Packages (accessed \today).; Java SE API docs (e.g., \texttt{Comparable}, \texttt{Runnable}) (accessed \today).}

\newcommand{\OOPSUnitIIISources}{Schildt, \textit{Java: The Complete Reference} (12e) (nested/inner classes, exceptions, threads, I/O).; Oracle Java Tutorials: Nested Classes + Exceptions + Concurrency + Basic I/O (accessed \today).; Java SE API docs (e.g., \texttt{Thread}, \texttt{AutoCloseable}) (accessed \today).}

\newcommand{\OOPSUnitIVSources}{Schildt, \textit{Java: The Complete Reference} (12e) (generics, lambdas, Swing, JDBC).; Bloch, \textit{Effective Java} (3e) (generics + lambdas).; Oracle Java Tutorials: Generics + Lambda Expressions + Swing + JDBC Basics (accessed \today).; Java SE API docs (e.g., \texttt{java.util.function}, \texttt{javax.swing}, \texttt{java.sql}) (accessed \today).}

\newcommand{\OOPSUnitVSources}{Schildt, \textit{Java: The Complete Reference} (12e) (Collections Framework + wrapper classes).; Oracle Java docs: Collections Framework overview (accessed \today).; Java SE API docs (\texttt{java.util} collections; wrapper classes in \texttt{java.lang}) (accessed \today).}

\newcommand{\OOPSUnitVISources}{Git SCM: \textit{Pro Git} (2e) (version control workflow), accessed \today.; GitHub Docs (repositories, issues, pull requests), accessed \today.; Oracle Java Tutorials: Swing + JDBC (accessed \today).; JUnit 5 User Guide (accessed \today).; Apache Maven Project: Getting Started (accessed \today).; Gradle User Manual: Getting Started (accessed \today).; UML class diagrams: UML 2.x overview/spec (accessed \today).}

\newcommand{\OOPSMiscWhyInterfacesSources}{Schildt, \textit{Java: The Complete Reference} (12e) (interfaces).; Bloch, \textit{Effective Java} (3e) (prefer interfaces where appropriate).; Oracle Java Tutorials: Interfaces (accessed \today).; Java SE API docs (e.g., \texttt{Comparable}, \texttt{Runnable}) (accessed \today).}



\title[OOPS with Java]{Object Oriented Programming (CSEG1044)}
\subtitle{Unit 04 -- Lecture 02: Generic Methods + Bounds + Wildcards}
\begin{document}

\begin{frame}[fragile]
  \titlepage
  \vspace{0.5em}
  \slidesources{\OOPSUnitIVSources}
\end{frame}

\begin{frame}{Quick Links}
  \centering
  \hyperlink{sec:overview}{\beamerbutton{Overview}}\hspace{0.6em}
  \hyperlink{sec:concepts}{\beamerbutton{Key Topics}}\hspace{0.6em}
  \hyperlink{sec:demo}{\beamerbutton{Demo}}\hspace{0.6em}
  \hyperlink{sec:summary}{\beamerbutton{Summary}}
  \slidesources{\OOPSUnitIVSources}
\end{frame}

\begin{frame}{Agenda}
  \tableofcontents
  \slidesources{\OOPSUnitIVSources}
\end{frame}

\section{Overview}
\label{sec:overview}

\begin{frame}{Lecture Snapshot}
\begin{itemize}[<+->]
  \item \textbf{Unit focus:} Generics, Lambdas, Swing, and JDBC
  \item \textbf{Course thread:} Use Generic Methods + Bounds + Wildcards to move Campus Resource Manager toward a real desktop application with UI, callbacks, and persistence.
\end{itemize}
  \slidesources{\OOPSUnitIVSources}
\end{frame}

\begin{frame}{Recall Before You Start}
\begin{itemize}[<+->]
  \item \textbf{Generic classes}: This lecture extends the first generics lecture from type declarations to more flexible operations. Main reference: \texttt{Unit 04 / Lecture 01 slides/notes}.
  \item \textbf{Inheritance}: Bounds and wildcards only make sense if subtype relationships are already familiar. Main reference: \texttt{Unit 02 / Lecture 04 slides/notes}.
\end{itemize}
  \slidesources{\OOPSUnitIVSources}
\end{frame}


\begin{frame}{Learning Outcomes}
  \begin{itemize}[<+->]
    \item Write a generic method and explain why it reduces duplication.
    \item Use bounds (\texttt{extends}) to restrict allowed types.
    \item Use wildcards (\texttt{? extends}, \texttt{? super}) with collections.
  \end{itemize}
  \slidesources{\OOPSUnitIVSources}
\end{frame}

\begin{frame}{Key Terms to Know}
\small
\begin{columns}[t]
\column{0.48\textwidth}
\begin{itemize}
  \item \textbf{Generic method}: method with its own type parameter, e.g., \texttt{static <T> T first(List<T> list)}.
  \item \textbf{Bound}: restriction on type parameter, e.g., \texttt{<T extends Number>}.
  \item \textbf{Wildcard}: unknown type parameter, e.g., \texttt{List<?>}.
\end{itemize}
\column{0.48\textwidth}
\begin{itemize}
  \item \textbf{Upper-bounded wildcard}: \texttt{? extends Type} (read/produce).
  \item \textbf{Lower-bounded wildcard}: \texttt{? super Type} (write/consume).
  \item \textbf{Invariance}: \texttt{List<Integer>} is not a subtype of \texttt{List<Number>}.
\end{itemize}
\end{columns}
  \slidesources{\OOPSUnitIVSources}
\end{frame}

\begin{frame}{Study Flow for This Lecture}
\begin{enumerate}[<+->]
  \item Refresh the prerequisite bridge before reading new ideas in isolation.
  \item Study the main build in this order: \textit{Generic methods} $\rightarrow$ \textit{Bounds (restrict allowed types)} $\rightarrow$ \textit{Wildcards: making collection parameters flexible}.
  \item Trace the worked example and connect each line back to one concept frame.
  \item Attempt the mini exercises before moving to the recap and exit check.
\end{enumerate}
  \slidesources{\OOPSUnitIVSources}
\end{frame}


\section{Key Topics}
\label{sec:concepts}

\begin{frame}{Unit 04: Generics, Lambdas, GUI Swing \& Database Connectivity}
  \begin{itemize}[<+->]
    \item Lecture focus: Generic Methods + Bounds + Wildcards
  \end{itemize}
  \slidesources{\OOPSUnitIVSources}
\end{frame}

\begin{frame}{Today: Roadmap}
  \begin{itemize}[<+->]
    \item Generic methods
    \item Bounded type parameters and wildcards (intro level)
    \item Using generics with collections (patterns that show up in real code)
  \end{itemize}
  \slidesources{\OOPSUnitIVSources}
\end{frame}

\begin{frame}{Generic method (idea)}
  \begin{itemize}[<+->]
    \item Make only the method generic (not necessarily the class).
    \item Syntax: \texttt{public static <T> void printArray(T[] a)}
  \end{itemize}
  \slidesources{\OOPSUnitIVSources}
\end{frame}

\begin{frame}{Bounds (restrict types)}
  \begin{itemize}[<+->]
    \item \texttt{<T extends Number>} limits \texttt{T} to numeric types.
    \item Common comparison pattern: \texttt{<T extends Comparable<T>>}
  \end{itemize}
  \slidesources{\OOPSUnitIVSources}
\end{frame}

\begin{frame}{Why wildcards? (invariance)}
  \begin{itemize}[<+->]
    \item \texttt{List<Integer>} is \textbf{not} a \texttt{List<Number>}.
    \item We need flexible method parameters for collections.
  \end{itemize}
  \slidesources{\OOPSUnitIVSources}
\end{frame}

\begin{frame}{Wildcards: \texttt{extends} vs \texttt{super}}
  \begin{itemize}[<+->]
    \item \texttt{? extends T}: use when list produces/read values of \texttt{T}.
    \item \texttt{? super T}: use when list consumes/you add values of \texttt{T}.
    \item PECS: Producer Extends, Consumer Super.
  \end{itemize}
  \slidesources{\OOPSUnitIVSources}
\end{frame}

\begin{frame}{Checkpoint and Reflection}
\begin{block}{Reflection prompt}
Where would \textit{Generic Methods + Bounds + Wildcards} matter most in Campus Resource Manager, and which design choice would you justify using this lecture's ideas?
\end{block}
\begin{block}{Quick self-check}
Which part needs one more example before you move on: \textit{Generic methods}, \textit{Bounds (restrict allowed types)}, the worked example, or the recap?
\end{block}
  \slidesources{\OOPSUnitIVSources}
\end{frame}

\section{Demo / Example}
\label{sec:demo}

\begin{frame}[fragile]{Example: generic method}
\begin{lstlisting}
public static <T> void printArray(T[] a) {
    for (T x : a) System.out.print(x + " ");
    System.out.println();
}
\end{lstlisting}
  \slidesources{\OOPSUnitIVSources}
\end{frame}

\begin{frame}[fragile]{Example: bounded \texttt{max}}
\begin{lstlisting}
public static <T extends Comparable<T>> T max(T a, T b) {
    return (a.compareTo(b) >= 0) ? a : b;
}
\end{lstlisting}
  \slidesources{\OOPSUnitIVSources}
\end{frame}

\begin{frame}[fragile]{Example: wildcard \texttt{sum}}
\begin{lstlisting}
import java.util.List;
public static double sum(List<? extends Number> list) {
    double s = 0;
    for (Number n : list) s += n.doubleValue();
    return s;
}
\end{lstlisting}
  \slidesources{\OOPSUnitIVSources}
\end{frame}

\begin{frame}{Mini Exercises (2--3 mins)}
  \begin{enumerate}[<+->]
    \item Why is \texttt{List<Integer>} not a \texttt{List<Number>}?
    \item Which wildcard for reading: \texttt{extends} or \texttt{super}?
    \item State the PECS rule.
  \end{enumerate}
  \slidesources{\OOPSUnitIVSources}
\end{frame}

\begin{frame}{Watch-outs}
\begin{itemize}[<+->]
  \item Using a wildcard when a named type parameter would be clearer.
  \item Overengineering beginner APIs with too many bounds.
  \item Assuming List<Child> is automatically a List<Parent>.
\end{itemize}
  \slidesources{\OOPSUnitIVSources}
\end{frame}

\section{Summary}
\label{sec:summary}

\begin{frame}{Key Takeaways}
  \begin{itemize}[<+->]
    \item Generic methods reduce duplication.
    \item Bounds restrict allowed types (safer APIs).
    \item Wildcards make collection parameters flexible (PECS).
  \end{itemize}
  \slidesources{\OOPSUnitIVSources}
\end{frame}

\begin{frame}{Checkpoint Questions}
  \begin{enumerate}[<+->]
    \item Write one example of a generic method signature.
    \item What does \texttt{<T extends Number>} restrict?
    \item When do you use \texttt{? super T}?
  \end{enumerate}
  \slidesources{\OOPSUnitIVSources}
\end{frame}

\begin{frame}{Next Study Steps}
\begin{itemize}[<+->]
  \item Revisit the recall references if any older idea still feels weak.
  \item Rework the lecture example without looking at the notes first.
  \item Attempt the self-study practice and then answer the exit questions in full sentences.
  \item Apply the idea once inside Campus Resource Manager to make the concept stick.
\end{itemize}
  \slidesources{\OOPSUnitIVSources}
\end{frame}

\begin{frame}{Next Lecture Preview}
  \textbf{Next:} Lambdas + Functional Interfaces
  \vspace{0.6em}
  \begin{itemize}[<+->]
    \item We will write short callback functions and use them with collections and UI events.
  \end{itemize}
  \slidesources{\OOPSUnitIVSources}
\end{frame}

\end{document}
